\chapter{INTRODUCTION}
\label{ch:intro}

\section{TWO WHEELED INVERTED PENDULUM}

The two-wheeled inverted pendulum (TWIP) is a mobile robot consisting of a body
mounted on a single wheel axis, with the center of gravity located above that
axis. The system is statically unstable and must be actively controlled to
remain upright, which makes it a convenient laboratory analogue for a broad
class of unstable, underactuated aerospace systems: it has more degrees of
freedom than control inputs, it has strongly coupled translational and
rotational dynamics, and its linearization is valid only in a neighborhood of
the operating point.

The TWIP is attractive as a testbed for three reasons. First, its dynamics are
well understood and can be derived from first principles, so a truth model is
available against which estimator performance can be assessed. Second, the
uncertainty in a physical build is real and not synthetic: mass properties are
difficult to measure accurately, the motor constants drift with temperature and
load, and the gear train exhibits deadzone and backlash. Third, the system is
genuinely underactuated, with a single motor voltage command driving both the
translational and rotational accelerations, so the uncertainties entering the
two acceleration equations cannot both be cancelled directly. The uncertainty is
therefore \emph{unmatched}, and this is the case of practical interest.

\section{CONTRIBUTION OF THIS WORK}

This work makes three contributions.

\begin{enumerate}[label=(\roman*),leftmargin=*]

\item \textbf{A discrete-time Modified State Observer with a complete stability
proof.} The MSO of Rajagopal et al.\ is extended to discrete time. Every digital
implementation is a discrete-time system, so a discrete-time result is required
before the technique can be claimed for digital flight control, rover autonomy,
or hardware-in-the-loop use. \rev{The proof given in
Chapter~\ref{ch:dmso} evaluates the Lyapunov first difference exactly. The
resulting admissibility conditions are $\sigmax(F-\Km)<1$ together with an upper
bound on the adaptation gain, both of which are satisfiable over an open set of
designs.}

\item \textbf{A command-filtered neural control law for unmatched uncertainty.}
Chapter~\ref{ch:control} develops an additive control signal that uses the tilt
angle as a virtual control for the position subsystem. \rev{The design is posed
as a four-step backstepping procedure with command filtering, which assigns
damping to every error coordinate, retains and cancels the coupling between the
position and attitude error subsystems, and makes each neural network's
approximation target a function of measured state and reference trajectory
alone.}

\item \textbf{Hardware validation.} The observer is implemented on a physical
TWIP robot\rev{, and the model of the robot is tested against its recorded
balancing runs (Section~\ref{sec:modelsim})}. Chapter~\ref{ch:platform} documents the platform and the practical
obstacles encountered; Appendix~\ref{app:validation} gives a staged validation
procedure for reproducing and extending the results.

\end{enumerate}

\section{ORGANIZATION}

Chapter~\ref{ch:lit} reviews the relevant literature on two-wheeled inverted
pendulum control and on modified state observers. Chapter~\ref{ch:dmso} develops
the DMSO and its stability proof. Chapter~\ref{ch:platform} describes the
physical robot. Chapter~\ref{ch:model} derives the TWIP equations of motion from
the DC motor model upward \rev{and tests them against the robot's recorded
data}. Chapter~\ref{ch:control} presents the LQR baseline
and the command-filtered neural controller. Chapter~\ref{ch:simulation}
describes the simulation environment, including the discretization, the noise
models, and the actuator nonlinearities. Chapter~\ref{ch:results} presents
results. Chapter~\ref{ch:conclusions} concludes.
Appendix~\ref{app:validation} gives implementation and test guidance;
Appendix~\ref{app:errata} records the specific technical differences between
this revision and revision~7.
